\documentclass[11pt,a4paper]{article}

\usepackage[utf8]{inputenc}
\usepackage[T1]{fontenc}
\usepackage[brazil]{babel}
\usepackage{geometry}
\usepackage{amsmath}
\usepackage{booktabs}
\usepackage{longtable}
\usepackage{array}
\usepackage{enumitem}
\usepackage{xcolor}
\usepackage{hyperref}

\geometry{
    left=2.5cm,
    right=2.5cm,
    top=2.5cm,
    bottom=2.5cm
}

\hypersetup{
    colorlinks=true,
    linkcolor=blue,
    urlcolor=blue,
    citecolor=blue,
    pdftitle={Especificacao do Estado Atual do Coprocessador Vetorial RISC-V},
    pdfauthor={Projeto autoral}
}

\setlist[itemize]{topsep=3pt,itemsep=2pt}
\setlist[enumerate]{topsep=3pt,itemsep=2pt}

\title{
    \textbf{Especificação do Estado Atual do Projeto}\\
    \large Modelo Executável de um Coprocessador Vetorial RISC-V
}
\author{Projeto autoral}
\date{\today}

\begin{document}

% \maketitle

\begin{titlepage}
    \centering

    \vspace*{0.22\textheight}

    {\Huge\bfseries UBIRA-V\par}
    \vspace{0.5cm}
    {\LARGE Especificacao de Projeto de um Coprocessador Vetorial RISC-V\par}

    \vspace{1.5cm}

    {\Large Lucas Farias Martins\par}

    \vfill

    {\large 12 de outubro de 2026\par}
\end{titlepage}

\clearpage

% =============================================================
%                _   _       _               _   
%               /_\ | |__ __| |_ _ _ __ _ __| |_ 
%              / _ \| '_ (_-<  _| '_/ _` / _|  _|
%             /_/ \_\_.__/__/\__|_| \__,_\__|\__|                                    
% =============================================================

\begin{abstract}
    Este documento registra o estado atual do projeto de um coprocessador
    vetorial inspirado na arquitetura RISC-V. O foco desta etapa é a
    especificação funcional executável em Python, anterior à migração para
    SystemC e à implementação RTL em SystemVerilog.
    
    A configuração física adotada é VLEN igual a 128 bits. O modelo está
    sendo generalizado para suportar elementos inteiros de 8, 16 e 32 bits,
    bem como grupos de registradores com LMUL igual a 1, 2, 4 ou 8. A
    especificação também descreve operações inteiras, cargas e stores
    contíguos, execução sequencial de programas, política de cauda e
    processamento de vetores maiores que VLMAX por strip-mining.
    
    O documento distingue decisões de projeto, comportamento pretendido,
    funcionalidades presentes no modelo e pendências identificadas durante
    a execução dos testes.
\end{abstract}
    
\clearpage

\tableofcontents

\clearpage

% =============================================================
%                      OBJETIVO E ESCOPO
% =============================================================

\section{Objetivo e escopo}

O objetivo do projeto é desenvolver um coprocessador vetorial para
pesquisa arquitetural, com foco futuro em uma implementação sintetizável
para ASIC.

O fluxo de desenvolvimento definido é:

\begin{enumerate}
    \item modelo funcional executável em Python;
    \item modelo temporal ou de arquitetura em SystemC;
    \item implementação RTL em SystemVerilog;
    \item simulação e verificação com Xcelium e SimVision;
    \item síntese e análise de área e temporização com Genus.
\end{enumerate}

Esta versão documenta o modelo funcional. Ela não especifica pipeline,
latência de hardware, protocolo de memória RTL, exceções arquiteturais
completas, interface com um processador escalar ou desempenho físico.

\clearpage

% =============================================================
%                ___         _   _             
%               / __| ___ __| |_(_)___ _ _  ___
%               \__ \/ -_) _|  _| / _ \ ' \(_-<
%               |___/\___\__|\__|_\___/_||_/__/
% =============================================================

\input{tex_src/01_status}
\input{tex_src/02_parameters}
\input{tex_src/03_architecture}
\input{tex_src/04_semantics}
\input{tex_src/05_software}
\input{tex_src/06_verification}
\input{tex_src/07_next_steps}
% \input{tex_src/08_revision_history}

\end{document}